\ncdlesson{19}{Udkantsdanmark eller storbydrømme?}{边缘乡村还是大城市之梦？}
{DU3 Modul 3（B1）｜对标 NCE 3 前半 \& Modul 3 \textit{Bolig, Geografi og Livsvalg}}
{\textbf{城乡对比说明文的段落衔接架构（平行对比法 vs 逐点对比法）} · 高阶对比衔接词与语序矩阵（\textit{hvorimod} 从句 vs \textit{derimod / til gengæld / på den anden side} 倒装）}
{DU3 Modul 3 的阅读理解与写作常常让你在“两种生活模式”之间做对比与权衡。写对比段落最怕“一段城市、下一句突然跳乡村、最后又跳回城市”。学完本课的两种黄金段落架构，你就能用清晰的对比连词写出让考官眼前一亮的 B1 结构化文章！}

% ---------------------------------------------------------------------
\begin{foerduleaeser}
带着下面两个问题阅读 Tekst A 和 Tekst B：
\begin{enumerate}
  \item \textit{Hvilke fordele og ulemper nævnes i Tekst A ved at bo i henholdsvis storbyen og Udkantsdanmark?}（在 Tekst A 中，分别提到了住在大城市和边缘乡村地区的哪些优势与劣势？）
  \item \textit{Hvad overvejer Jonas og Sara i deres chat i Tekst B, og hvilken løsning vælger de til sidst?}（在 Tekst B 的聊天中，乔纳斯和萨拉在纠结什么？他们最后倾向哪种方案？）
\end{enumerate}
\end{foerduleaeser}

% ---------------------------------------------------------------------
\ncdsubsec{二、Tekst A \& Tekst B（社会深读与 Modul 3 真题场景）}

\subsection*{Tekst A：\textit{Udkantsdanmark eller storbydrømme?}（社会深读 · 约 250 词）}

\begin{tipbox}{语篇说明}
本篇采用\textbf{“逐点对比法”（Punkt-for-punkt-metoden）}依次从\ding{172}住房成本、\ding{173}就业通勤、\ding{174}生活节奏三个维度展开城乡对比。请重点观察\textbf{\color{NCDNavy}深蓝色标注的句首对比过渡短语}与\textbf{\color{NCDRed}深红色标注的从句/副词对比连接词}。
\end{tipbox}

Selvom Danmark geografisk er et lille land, er forskellen mellem de store byer og de såkaldte yderområder --- ofte kaldet ”Udkantsdanmark” --- blevet tydeligere i de seneste år. Mange unge drømmer om København eller Aarhus på grund af studier og et aktivt byliv, mens børnefamilier og pensionister oftere kigger mod mindre byer og landdistrikter, hvor hverdagen føles mere overskuelig.

\da{På den ene side} er huslejen og boligpriserne i storbyerne steget så kraftigt, at det kan være svært for almindelige lønmodtagere at købe en større bolig. \da{På den anden side} får man langt mere plads for pengene i provinsen eller på små øer, \dabold{hvorimod} udvalget af små, centrale lejligheder her kan være begrænset. Mange flytter netop ud for at få en privat have til børnene.

Når det gælder arbejde og uddannelse, ligger de store universiteter og internationale virksomheder stadig tæt i storbyerne, \da{mens} man i landdistrikterne \dabold{derimod} ofte må acceptere en længere pendlertid og færre offentlige busser. I de senere år har mere fleksibelt hjemmearbejde dog gjort det lettere for nogle pendlere at bo længere væk fra kontoret, uden at de mister deres karrieremuligheder.

\da{Sammenlignet med} storbyens travle tempo tilbyder Udkantsdanmark frisk luft, tæt natur og et stærkt lokalsamfund, hvor naboerne kender hinanden. \dabold{Til gengæld} savner nogle tilflyttere storbyens brede kulturliv, mange restauranter og internationale miljø. Valget mellem by og land handler derfor sjældent om, hvad der er ”bedst” for alle, men snarere om, hvilken livsfase man befinder sig i.

\begin{transbox}
\small
\textbf{\color{NCDTeal}[Tekst A 中文精译]}\\
尽管丹麦在地理上是个小国，但大城市与所谓外围地区——常被称为“边缘丹麦（Udkantsdanmark）”——之间的差异近年来变得越来越明显。许多年轻人因为求学和丰富活跃的城市生活而向往哥本哈根或奥胡斯，而有孩子的家庭和退休人员则更常把目光投向日常生活显得更从容可控的中小城镇和乡村地区。

一方面，大城市的房租和房价上涨得如此之快，以至于普通工薪阶层很难买到一套较大的住房。另一方面，在地方城镇或小岛上，同样的钱可以换来大得多的居住空间，不过当地市中心小户型的选择相对有限。许多家庭正是为了给孩子们换取一个私人花园才搬出城市。

至于工作和教育，大型高校和跨国企业依然密集分布在大城市，而在乡村地区则往往不得不接受更长的通勤时间和更少的公共交通班次。不过近年来，更灵活的居家办公方式让一些通勤者即使住得离办公室更远，也不会失去职业发展机会。

与大城市忙碌的节奏相比，丹麦乡村地区提供了清新的空气、亲近的自然以及邻里彼此相熟的紧密地方社区。不过相应地，也有一些新搬来的人会怀念大城市丰富的文化生活、众多餐馆和国际化氛围。因此，选择城市还是乡村，很少取决于哪种对所有人来说“最好”，而更多取决于一个人正处于哪个人生阶段。
\end{transbox}

\subsection*{Tekst B：\textit{Flytte ud af storbyen?}（对标 DU3 Modul 3 \textit{Læsning Opgave 2A}：对话气泡补全）}

\textbf{阅读任务（Opgave 2A 模拟）：}\textit{Situation: Jonas skriver til sin kæreste, Sara, som er på arbejde, om deres boligplaner. Læs Jonas' beskeder, og find de svar fra Sara (A--F), der passer til hullerne (1)--(3). Der er tre svar, du ikke skal bruge. Se eksemplet (0).}

\begin{emailbox}{Læsning Opgave 2A｜Chat mellem Jonas og Sara}
\textbf{Jonas (0):} Hej Sara! Har du haft tid til at kigge på det lille hus i Silkeborg, som jeg sendte i går?\\
\textbf{Sara (0):} \textbf{\underline{\ Z\ }} \textit{Ja, det har jeg lige set i frokostpausen. Det ser virkelig hyggeligt ud med den store have!}

\vspace{0.3em}
\textbf{Jonas:} Ikke sandt? Men jeg er stadig lidt nervøs for afstanden. Jeg skal jo pendle næsten en time til Aarhus hver dag.\\
\textbf{Sara (1):} \underline{\hspace{1.2cm}}

\vspace{0.3em}
\textbf{Jonas:} Det har du ret i. Og min chef sagde faktisk i morges, at jeg godt må arbejde hjemme to dage om ugen.\\
\textbf{Sara (2):} \underline{\hspace{1.2cm}}

\vspace{0.3em}
\textbf{Jonas:} Ja, bestemt! På den anden side vil vi jo savne vores venner og caféerne i midtbyen lidt.\\
\textbf{Sara (3):} \underline{\hspace{1.2cm}}

\vspace{0.3em}
\textbf{Jonas:} Super idé! Jeg skriver til ejendomsmægleren nu og bestiller en fremvisning på lørdag.

\vspace{0.4em}
\hrule
\vspace{0.3em}
\small
\textbf{Saras mulige svar（A--F，6 选 3，其中 3 项多余）：}
\begin{itemize}[label={}, leftmargin=0.8em]
  \item \textbf{A.} Til gengæld ligger mine forældres hus kun ti minutter derfra, og vores venner kan jo altid komme på besøg i weekenden!
  \item \textbf{B.} Ej, hvor er det ærgerligt, at du slet ikke kan arbejde hjemmefra i dit nye job.
  \item \textbf{C.} Det er rigtigt, men til gengæld er huslejen jo næsten halvt så dyr som i vores lille lejlighed her.
  \item \textbf{D.} Derfor synes jeg hellere, vi skal købe en dyrere lejlighed lige midt i Aarhus.
  \item \textbf{E.} Ej, hvor godt! Det gør jo en kæmpe forskel for pendlertiden. Skal vi så tage ud at se huset i weekenden?
  \item \textbf{F.} Sammenlignet med Silkeborg er vores have her i Aarhus desværre alt for stor at passe.
\end{itemize}
\end{emailbox}

\begin{tipbox}{本课发音与对比连接词读音要点（Udtale \& Reduktioner）}
\begin{itemize}
  \item \textbf{对比连词与副词的读音及不发音字母}：
    \begin{itemize}
      \item \textit{hvorimod} \ipa{vɒːiˈmoðˀ}（词首 \textit{h} 不发音！重音在最后的 \textit{-mod} 上，软 d + \textit{stød}）
      \item \textit{derimod} \ipa{ˈdɛɐ̯iˌmoðˀ}（重音在前面的 \textit{der-} 上，刚好与 \textit{hvorimod} 相反！）
      \item \textit{til gengæld} \ipa{tˢe ˈɡɛŋˌɡɛlˀ}（\textit{til} 弱读 \ipa{tˢe}；\textit{gæld} 中的 \textit{d} 不发音！）
      \item \textit{sammenlignet} \ipa{ˈsɑmənˌliːˀnəð}（\textit{g} 不发音，词尾 \textit{-et} 读轻弱的 \ipa{əð}）
    \end{itemize}
  \item \textbf{城乡生活高频词}：\textit{Udkantsdanmark} \ipa{ˈuðkʰanˀdsˌdanmɑɡ}（\textit{ud} 读软 d，\textit{kants} 带 \textit{stød}）；\textit{storbydrømme} \ipa{ˈsdoɐ̯byˌdʁœmə}；\textit{et lokalsamfund} \ipa{loˈkʰæːˀlˌsɑmfɔnˀ}（词尾 \textit{d} 不发音）；\textit{en husleje} \ipa{ˈhuːsˌlɑɪ̯ə}。
\end{itemize}
\end{tipbox}

\begin{transbox}
\small
\textbf{\color{NCDTeal}[Tekst B 中文对照]}\\
\textbf{聊天对话（搬离大城市吗？）}：
\textbf{乔纳斯 (0)}：嗨萨拉！你有没有抽空看我昨天发给你的锡尔克堡那栋小房子？\\
\textbf{萨拉 (0/Z)}：看了，我午休时刚看！带着大花园看起来真的好温馨！\\
\textbf{乔纳斯}：是吧？但我还是有点担心距离问题。我每天得往奥胡斯通勤差不多一小时呢。\\
\textbf{萨拉 (1)}：\textbf{[空 1]} \quad \textbf{乔纳斯}：你说得对。而且我老板今天早上还真说了，我每周可以居家办公两天。\\
\textbf{萨拉 (2)}：\textbf{[空 2]} \quad \textbf{乔纳斯}：那是当然！不过另一方面，我们也会有点想念市中心的朋友和咖啡馆。\\
\textbf{萨拉 (3)}：\textbf{[空 3]} \quad \textbf{乔纳斯}：好主意！我现在就给房产中介写信，预约周六去看房。
\end{transbox}

% ---------------------------------------------------------------------
\ncdsubsec{三、Gloser（核心生词表）}

\begin{tipbox}{词表重点}
注意 \textit{at tiltrække}（吸引）的不规则过去时是 \dabold{tiltrak}（不是 \textit{tiltræk}），过去分词是 \dabold{har tiltrukket}。另外，\textit{at trives}（生活舒心/茁壮适应）是一个\textbf{异态动词（Deponent verbum）}——外形永远带 \textit{-s}，但意思完全是主动的（如 \textit{Vi trives godt på landet}）。
\end{tipbox}

\begin{glosertable}
名词 & et landdistrikt & \textit{(landdistriktet, landdistrikter, -erne)} & 乡村地区，农村区域 \\
名词 & en storby & \textit{(storbyen, storbyer, storbyerne)} & 大城市，都市 \\
名词 & et lokalsamfund & \textit{(lokalsamfundet, lokalsamfund, -ene)} & 地方社区，本地社群 \\
名词 & en husleje & \textit{(huslejen, huslejer, huslejerne)} & 房租 \\
名词 & en boligpris & \textit{(boligprisen, boligpriser, -erne)} & 房价 \\
名词 & en fordel & \textit{(fordelen, fordele, fordelene)} & 优点，优势 \\
名词 & en ulempe & \textit{(ulempen, ulemper, ulemperne)} & 缺点，劣势 \\
名词 & et kulturliv & \textit{(kulturlivet, ingen flertal)} & 文化生活 \\
名词 & en pendlertid & \textit{(pendlertiden, pendlertider, -erne)} & 通勤时间 \\
名词 & en livskvalitet & \textit{(livskvaliteten, ingen flertal)} & 生活质量 \\
\midrule
动词 [har] & at tiltrække & \textit{(tiltrækker, tiltrak, har tiltrukket)} & 吸引 \\
动词 \dabold{[er]} & at flytte & \textit{(flytter, flyttede, \dabold{er flyttet})} & 搬家，迁居 \\
动词 [har] & at sammenligne (med) & \textit{(sammenligner, sammenlignede, har sammenlignet)} & 把……与……作比较 \\
动词 [har] & at overveje & \textit{(overvejer, overvejede, har overvejet)} & 考虑，斟酌 \\
异态动词 & at trives & \textit{(trives, trivedes, har trivedes)} & 舒心生活，适应良好 \\
\midrule
形容词 & rolig & \textit{(roligt, rolige)} & 安静的，平缓的 \\
形容词 & travl & \textit{(travlt, travle)} & 忙碌的，繁忙的 \\
\end{glosertable}

% ---------------------------------------------------------------------
\ncdsubsec{四、Grammatik \& Faste vendinger（核心语法与地道表达精讲）}

\begin{grammatikbox}{语法点 1：城乡对比说明文的两大黄金段落架构}
在 DU3 写作（如邮件建议、社会议题短文）中对比 A 与 B 两个选项时，切忌想到哪写到哪，请根据字数选择下面两种标准架构之一：

\vspace{0.3em}
\renewcommand{\arraystretch}{1.24}
{\small
\begin{tabularx}{\linewidth}{@{} >{\bfseries\raggedright\arraybackslash}p{3.0cm} >{\raggedright\arraybackslash}p{5.5cm} >{\raggedright\arraybackslash}X @{}}
\toprule
\textbf{段落架构} & \textbf{段落排布模板} & \textbf{适用场景与段落路标} \\
\midrule
\textbf{1. 平行对比法}\newline （\textit{Blokmetoden}） &
第 1 段：引言提出选择\newline
第 2 段：\textbf{集中写完方案 \textit{A}（如大城市）}的优缺点\newline
第 3 段：\textbf{集中写完方案 \textit{B}（如乡村）}的优缺点\newline
第 4 段：权衡总结 &
\textbf{最适合 100--150 词短文与回信}。\newline
第 3 段开头强转场：\textit{På den anden side / Hvis man derimod vælger at bo på landet...} \\
\midrule
\textbf{2. 逐点对比法}\newline （\textit{Punkt-for-punkt}） &
第 1 段：引言提出对比背景\newline
第 2 段：\textbf{维度一（如房价/空间）}：A vs B\newline
第 3 段：\textbf{维度二（如工作/交通）}：A vs B\newline
第 4 段：\textbf{维度三（如社交/结论）}：A vs B &
\textbf{最适合 200--300 词长说明文}（见本课 Tekst A）。\newline
每段开头亮维度：\textit{Når det gælder... / Hvad angår... / Sammenlignet med...} \\
\bottomrule
\end{tabularx}
}
\end{grammatikbox}

\begin{grammatikbox}{语法点 2：五大高阶对比衔接词的“词性与语序矩阵”}
同样表示“相反 / 而 / 另一方面”，它们放进句子里触发的语序规则完全不同！务必分清\textbf{从属连词}和\textbf{副词/介词短语}：

\vspace{0.3em}
\renewcommand{\arraystretch}{1.24}
{\small
\begin{tabularx}{\linewidth}{@{} >{\bfseries\raggedright\arraybackslash}p{2.9cm} >{\raggedright\arraybackslash}p{3.5cm} >{\raggedright\arraybackslash}X @{}}
\toprule
\textbf{衔接词} & \textbf{词性与核心语序} & \textbf{标准句式示范（注意加粗词位置）} \\
\midrule
\da{hvorimod} & \textbf{从属连词}（“而，反之”）\newline 引导\textbf{从句：不倒装}，\textit{ikke} 放动词前 & \textit{I byen er huslejen høj, \da{hvorimod} man på landet \dabold{ikke betaler} så meget.} \\
\midrule
\da{mens ... derimod} & \textit{mens} 引导从句，\textit{derimod} 插在从句主语后（或句中） & \textit{Byen er travl, \da{mens} livet på landet \dabold{derimod er} meget roligere.} \\
\midrule
\da{Derimod} /\newline \da{Til gengæld} & \textbf{主句副词}\newline 放句首\textbf{必须 V2 倒装}（V + S） & \textit{\da{Til gengæld} \dabold{får man} mere plads.}\newline \textit{\da{Derimod} \dabold{er der} færre busser på landet.} \\
\midrule
\da{På den ene side..., på den anden side...} & \textbf{成对副词短语}\newline 放句首\textbf{两句都要 V2 倒装} & \textit{\da{På den ene side} \dabold{har byen} mange job, \da{på den anden side} \dabold{koster} boligerne en formue.} \\
\midrule
\da{Sammenlignet med + 名词} & \textbf{分词介词短语}\newline 放句首\textbf{主句必须 V2 倒装} & \textit{\da{Sammenlignet med København} \dabold{er} Silkeborg meget billigere.} \\
\bottomrule
\end{tabularx}
}
\end{grammatikbox}

\begin{grammatikbox}{英丹对比排雷 \& 地道搭配（Faste vendinger）}
\begin{enumerate}
  \item \textbf{别把 \textit{On the other hand} 的无倒装带进丹麦语！}：英语写 \textit{On the other hand, you get more space}；丹麦语 \textit{På den anden side} 占第 1 位，动词必须第 2 位：\ding{55}\ \textit{På den anden side man får mere plads} $\rightarrow$ \ding{51}\ \textit{På den anden side \dabold{får man} mere plads}！
  \item \textbf{严格区分 \textit{hvorimod} 与 \textit{derimod}}：\textit{hvorimod} 是连接两半句的\textbf{从属连词}（不单独领起独立句，后接主语 + 动词）；\textit{derimod} 是\textbf{副词}（可放独立句首触发 V2 倒装，或跟在 \textit{mens / og} 后用）。
  \item \textbf{\textit{til gengæld} 的地道用法}：意为“作为回报/相对应的是（有得也有失，或坏事里的弥补优点）”，极度高频！例如房贵但离朋友近，或通勤远但房子大，都要用 \textit{til gengæld}。
\end{enumerate}

\textbf{本课高频地道搭配（Faste vendinger）：}
\begin{itemize}
  \item \da{at få mere plads for pengene}（花同样的钱获得更大的空间）/ \da{at bo tæt på naturen}（住得亲近自然）
  \item \da{at have en høj livskvalitet}（拥有高质量生活）/ \da{at trives i lokalsamfundet}（在本地社区生活得融洽）
  \item \da{Når det gælder...} / \da{Hvad angår...}（谈到……方面——逐点对比法段落起手式！）
\end{itemize}
\end{grammatikbox}

% ---------------------------------------------------------------------
\ncdsubsec{五、Mønsterøvelser（句型操练）}

\textbf{练习 A：语序与对比衔接词填空（\textit{hvorimod, derimod, til gengæld, på den ene side, på den anden side, sammenlignet med}）}
\begin{enumerate}
  \item \underline{\hspace{3.2cm}} København er en lille by som Vejle meget roligere og billigere.
  \item I storbyen har man mange caféer og museer, \underline{\hspace{2.4cm}} man i landdistrikterne ofte har naturen lige uden for døren.
  \item \underline{\hspace{2.8cm}} er det nemt at finde arbejde i byen, \underline{\hspace{2.8cm}} kan huslejen være alt for høj.
  \item Der er langt til den nærmeste biograf fra vores nye hus. \underline{\hspace{2.4cm}} har vi nu en stor have og søde naboer!
  \item Mange unge flytter til København for at studere, mens børnefamilier \underline{\hspace{2.0cm}} ofte søger ud mod grønne forstæder.
  \item 找出并改正 1 处语序错误：\textit{På den anden side man skal køre i bil næsten hver dag på landet.} $\rightarrow$
  \item 找出并改正 1 处从句语序错误：\textit{Byen har mange koncerter, hvorimod på landet har man ikke så mange tilbud.} $\rightarrow$
\end{enumerate}

\textbf{练习 B：段落衔接“建筑手术”——按要求重组合并句子}
\begin{enumerate}
  \item 【逐点对比法｜住房维度】用 \textit{På den ene side..., på den anden side...} 合并以下两句：\newline
        \textit{En treværelses lejlighed i København koster ofte meget om måneden.} + \textit{Man kan leje et helt hus med have i Nordjylland for samme pris.} $\rightarrow$
  \item 【逐点对比法｜交通维度】用 \textit{hvorimod} 合并（注意从句中 \textit{ikke} 的位置！）：\newline
        \textit{I storbyen går metroen hvert femte minut.} + \textit{Man kan ikke altid stole på hyppige busser i små landsbyer.} $\rightarrow$
  \item 【平行对比法｜转场起句】用 \textit{Sammenlignet med} 触发 V2 倒装合并：\newline
        \textit{Vi ser på det travle byliv.} + \textit{Hverdagen i et lille lokalsamfund føles meget mere afslappet.} $\rightarrow$
  \item 【弥补转折】用 \textit{Til gengæld} 改写并连成通顺两句：\newline
        \textit{Det tager 40 minutter at pendle med tog. Vi bor kun fem minutter fra skoven og stranden.} $\rightarrow$
\end{enumerate}

\textbf{练习 C：中译丹（重点演练城乡对比连词与 V2 倒装）}
\begin{enumerate}
  \item 与大城市相比，丹麦乡村的房价通常低得多。
  \item 一方面大学和工作机会都在城市里，另一方面大城市的房租非常高。
  \item 许多年轻人喜欢热闹的城市生活，而家庭则反而更喜欢安静的环境。（用 \textit{mens ... derimod}）
  \item 通勤时间有点长；不过相应地，我们在小城里拥有很高的生活质量。（用 \textit{Til gengæld}）
  \item 在乡村邻里关系很紧密，而在大城市里人们有时根本不认识自己的邻居。（用 \textit{hvorimod} + \textit{ikke}）
\end{enumerate}

% ---------------------------------------------------------------------
\ncdsubsec{六、DU3 Modul 3 实战微任务与参考答案}

\begin{emailbox}{Skrivning Opgave 2 实战模拟｜给朋友的城乡选择建议信}
\textbf{Situation}: Din ven, Mikkel, er snart færdig med sin uddannelse. Han overvejer, om han skal blive boende i København eller flytte til en mindre by i Jylland for at få sit første job. Han har skrevet til dig for at høre din mening.\\
\textbf{Opgave}: Skriv et svar til Mikkel på \textbf{115--145 ord}. Sammenlign de to muligheder ved enten at bruge \textit{blokmetoden} (平行对比法) eller \textit{punkt-for-punkt-metoden} (逐点对比法), og giv ham et godt råd til sidst. Brug mindst tre kontrastudtryk fra denne lektion.
\end{emailbox}

\begin{answerbox}
\textbf{Før du læser 参考回答}
\begin{enumerate}
  \item \textit{Storbyen har universiteter, virksomheder og et bredt kulturliv, men høje boligpriser og et travlt tempo. Udkantsdanmark tilbyder mere plads for pengene, frisk luft, natur og et stærkt lokalsamfund, men kræver ofte længere pendlertid og færre busser.}
  \item \textit{De overvejer at flytte fra Aarhus til et lille hus i Silkeborg. Da Jonas må arbejde hjemme to dage om ugen, og huslejen er meget lavere, beslutter de at bestille en fremvisning af huset på lørdag.}
\end{enumerate}

\textbf{Tekst B 阅读答案（Opgave 2A）}
\begin{enumerate}
  \item[(1)] \textbf{C} —— Jonas 担心通勤一小时，Sara 用 \textit{Det er rigtigt, men til gengæld er huslejen jo næsten halvt så dyr...} 回应并提出补偿优势，自然引出 Jonas 下句 \textit{Det har du ret i.}
  \item[(2)] \textbf{E} —— Jonas 说老板允许每周居家办公两天，Sara 惊喜回应 \textit{Ej, hvor godt! Det gør jo en kæmpe forskel...}（B 项语义相反）。
  \item[(3)] \textbf{A} —— Jonas 担心想念市中心的朋友和咖啡馆，Sara 用 \textit{Til gengæld ligger mine forældres hus kun ti minutter derfra, og vores venner kan jo altid komme på besøg...} 宽慰并引出“Super idé!”。
\end{enumerate}

\textbf{练习 A 答案}
\begin{enumerate}
  \item \da{Sammenlignet med}（后接名词短语，注意句中 \textit{er en lille by...} 倒装）
  \item \da{hvorimod}（从属连词，从句中主语 \textit{man} 在动词 \textit{har} 前）
  \item \da{På den ene side}；\da{på den anden side}（两句均已 V2 倒装）
  \item \da{Til gengæld}（独立句首触发倒装 \textit{har vi}）
  \item \da{derimod}（位于 \textit{mens} 从句主语之后）
  \item 缺少 V2 倒装 $\rightarrow$ \textit{På den anden side \dabold{skal man} køre i bil næsten hver dag på landet.}
  \item 从句内部不能用主句倒装！$\rightarrow$ \textit{Byen har mange koncerter, \da{hvorimod} man på landet \dabold{ikke har} så mange tilbud.}
\end{enumerate}

\textbf{练习 B 答案}
\begin{enumerate}
  \item \textit{\da{På den ene side} koster en treværelses lejlighed i København ofte meget om måneden, \da{på den anden side} kan man leje et helt hus med have i Nordjylland for samme pris.}
  \item \textit{I storbyen går metroen hvert femte minut, \da{hvorimod} man i små landsbyer \dabold{ikke} altid kan stole på hyppige busser.}
  \item \textit{\da{Sammenlignet med} det travle byliv \dabold{føles} hverdagen i et lille lokalsamfund meget mere afslappet.}
  \item \textit{Det tager 40 minutter at pendle med tog. \da{Til gengæld} \dabold{bor vi} kun fem minutter fra skoven og stranden.}
\end{enumerate}

\textbf{练习 C 答案}
\begin{enumerate}
  \item \textit{Sammenlignet med storbyerne er boligpriserne i landdistrikterne som regel meget lavere.}
  \item \textit{På den ene side ligger universiteterne og mange arbejdspladser i byerne, på den anden side er huslejen i storbyen meget høj.}
  \item \textit{Mange unge kan lide det travle byliv, mens familier derimod foretrækker rolige omgivelser.}
  \item \textit{Pendlertiden er lidt lang. Til gengæld har vi en høj livskvalitet i den lille by.}
  \item \textit{På landet kender naboerne hinanden godt, hvorimod man i storbyen nogle gange slet ikke kender sine naboer.}
\end{enumerate}

\textbf{Modul 3 城乡建议高分回信（139 词 · 采用平行对比法）}\\
\textit{Til: Mikkel \quad|\quad Emne: Sv: Storby eller mindre by efter studiet?}

\textit{Hej Mikkel}

\textit{Stort tillykke med, at du snart er færdig med din uddannelse! Jeg kan godt forstå, at du overvejer begge muligheder.}

\textit{\da{På den ene side} er det nemmere at finde det første spændende job i København, og du har allerede et stort netværk og et rigt kulturliv lige uden for døren. \da{Sammenlignet med} mindre byer er huslejen dog meget høj, og det kan være svært at spare penge op.}

\textit{\da{På den anden side} får du i en mindre jysk by meget mere plads for pengene og et roligt tempo tæt på naturen. \da{Til gengæld} skal du måske regne med at bruge en bil, \da{hvorimod} du i København bare kan cykle overalt.}

\textit{Hvis du får et godt jobtilbud i Jylland, synes jeg, at du skal prøve det --- du kan jo altid flytte tilbage senere!}

\textit{Mange hilsner\\
\mbox{[Dit navn]}}

\small\textbf{考官点评}：全信结构严格按照“平行对比法”先集中论述哥本哈根、再用 \textit{På den anden side} 切换至日德兰小城，中间精准穿插了 \textit{Sammenlignet med}, \textit{Til gengæld}（V2 倒装）和 \textit{hvorimod}（从句语序），篇幅 139 词（远超 Modul 3 的 90 词底线），论证自然有力。
\end{answerbox}
